\documentclass[11pt,a4paper]{article}
\usepackage{amsmath,amssymb,amsfonts}
\usepackage{bm}

\title{Analytical Distance Calculations between von Mises and Wrapped Cauchy Distributions}
\author{Naoki Otani}
\date{\today}

\begin{document}
\maketitle

\section{Introduction}
This document summarizes the mathematical details of the analytical distance calculations (KL divergence and 1-Wasserstein distance) between the von Mises distribution and the wrapped Cauchy distribution on the circle $S^1$.

\section{Distributions}
Let $P$ be the von Mises distribution with mean direction $\mu_P \in [-\pi, \pi)$ and concentration parameter $\kappa > 0$:
\begin{equation}
p(\theta) = \frac{1}{2\pi I_0(\kappa)} e^{\kappa \cos(\theta - \mu_P)}
\end{equation}
where $I_n(\kappa)$ is the modified Bessel function of the first kind of order $n$.

Let $Q$ be the wrapped Cauchy distribution with mean direction $\mu_Q \in [-\pi, \pi)$ and concentration parameter $\rho \in [0, 1)$:
\begin{equation}
q(\theta) = \frac{1 - \rho^2}{2\pi (1 + \rho^2 - 2\rho \cos(\theta - \mu_Q))}
\end{equation}

\section{Kullback-Leibler Divergence}

\subsection{Divergence from von Mises to Wrapped Cauchy}
The KL divergence $D_{\text{KL}}(P \parallel Q)$ is defined as:
\begin{equation}
D_{\text{KL}}(P \parallel Q) = \int_0^{2\pi} p(\theta) \log \left( \frac{p(\theta)}{q(\theta)} \right) d\theta
\end{equation}
Using the Fourier series expansion:
\begin{equation}
\log (1 + \rho^2 - 2\rho \cos x) = -2 \sum_{n=1}^\infty \frac{\rho^n}{n} \cos(nx)
\end{equation}
and the trigonometric moments of the von Mises distribution:
\begin{equation}
E_P[\cos(n(\theta - \mu_Q))] = \frac{I_n(\kappa)}{I_0(\kappa)} \cos(n(\mu_P - \mu_Q))
\end{equation}
we obtain the fast-converging series representation:
\begin{equation}
D_{\text{KL}}(P \parallel Q) = \kappa \frac{I_1(\kappa)}{I_0(\kappa)} - \log I_0(\kappa) - \log(1 - \rho^2) - 2 \sum_{n=1}^\infty \frac{\rho^n I_n(\kappa)}{n I_0(\kappa)} \cos(n(\mu_P - \mu_Q))
\end{equation}

\subsection{Divergence from Wrapped Cauchy to von Mises}
The KL divergence $D_{\text{KL}}(Q \parallel P)$ can be calculated in closed-form. Using $E_Q[\cos(\theta - \mu_P)] = \rho \cos(\mu_Q - \mu_P)$ and $E_Q[\log(1 + \rho^2 - 2\rho \cos(\theta - \mu_Q))] = 2 \log(1 - \rho^2)$, we get:
\begin{equation}
D_{\text{KL}}(Q \parallel P) = \log I_0(\kappa) - \log(1 - \rho^2) - \kappa \rho \cos(\mu_Q - \mu_P)
\end{equation}

\section{1-Wasserstein Distance (Aligned Means)}
When the mean directions are aligned ($\mu_P = \mu_Q$) and one distribution is more concentrated than the other, the 1-Wasserstein distance $W_1(P, Q)$ on the circle reduces to the difference of their mean absolute deviations on $[-\pi, \pi]$:
\begin{equation}
W_1(P, Q) = | E_P[|\theta|] - E_Q[|\theta|] |
\end{equation}

\subsection{Derivation and Proof}
Without loss of generality, assume the aligned mean direction is $\mu_P = \mu_Q = 0$ and the distributions are defined on $[-\pi, \pi]$. Since both the von Mises and wrapped Cauchy distributions are symmetric around $\theta = 0$, their cumulative distribution functions (CDFs) satisfy $F_P(-\theta) = 1 - F_P(\theta)$ and $F_Q(-\theta) = 1 - F_Q(\theta)$ for $\theta \in [0, \pi]$.

The 1-Wasserstein distance on the circle $S^1$ is given by:
\begin{equation}
W_1(P, Q) = \min_{\alpha \in \mathbb{R}} \int_{-\pi}^\pi |F_P(\theta) - F_Q(\theta) - \alpha| d\theta
\end{equation}
Let $D(\theta) = F_P(\theta) - F_Q(\theta)$. By the symmetry of the CDFs, we have:
\begin{equation}
D(-\theta) = (1 - F_P(\theta)) - (1 - F_Q(\theta)) = -D(\theta)
\end{equation}
which implies that the difference function $D(\theta)$ is odd. The optimal shift $\alpha^*$ that minimizes the integral is the median of $D(\theta)$ over $[-\pi, \pi]$. Since $D(\theta)$ is a continuous odd function on the symmetric interval $[-\pi, \pi]$, its median is exactly $0$. Therefore, $\alpha^* = 0$, and the distance simplifies to:
\begin{equation}
W_1(P, Q) = \int_{-\pi}^\pi |F_P(\theta) - F_Q(\theta)| d\theta
\end{equation}
Since $|D(\theta)|$ is an even function, we can restrict the integration to $[0, \pi]$:
\begin{equation}
W_1(P, Q) = 2 \int_0^\pi |F_P(\theta) - F_Q(\theta)| d\theta
\end{equation}

Now, we examine the condition under which the CDF difference $D(\theta) = F_P(\theta) - F_Q(\theta)$ does not change sign (i.e., has 0 crossings) on the open interval $(0, \pi)$. For symmetric, unimodal circular distributions centered at $0$, the derivative $D'(\theta) = p(\theta) - q(\theta)$ can change sign at most twice in $(0, \pi)$. Since $D(0) = D(\pi) = 0$, the derivative must change sign at least once. 

A necessary and sufficient condition for $D(\theta)$ to have no crossings in $(0, \pi)$ is:
\begin{equation}
(p(0) - q(0))(p(\pi) - q(\pi)) \le 0
\end{equation}
Indeed:
\begin{itemize}
    \item If $(p(0) - q(0))(p(\pi) - q(\pi)) \le 0$, the derivative $D'(\theta)$ changes sign exactly once in $(0, \pi)$. Thus, $D(\theta)$ is unimodal and remains strictly positive (if $p(0) > q(0)$) or strictly negative (if $p(0) < q(0)$) on $(0, \pi)$, meaning it has 0 crossings.
    \item If $(p(0) - q(0))(p(\pi) - q(\pi)) > 0$, the derivative $D'(\theta)$ changes sign exactly twice in $(0, \pi)$. By Rolle's theorem and the Intermediate Value Theorem, this implies that $D(\theta)$ must cross $0$ exactly once in $(0, \pi)$.
\end{itemize}
Therefore, if the condition $(p(0) - q(0))(p(\pi) - q(\pi)) \le 0$ is satisfied, $F_P(\theta) - F_Q(\theta)$ does not change sign on $(0, \pi)$. Assuming without loss of generality that $F_P(\theta) \geq F_Q(\theta)$ for all $\theta \in [0, \pi]$, we can remove the absolute value:
\begin{equation}
W_1(P, Q) = 2 \int_0^\pi (F_P(\theta) - F_Q(\theta)) d\theta
\end{equation}

Next, we establish the relationship between this integral and the mean absolute deviation (MAD). The MAD of a symmetric distribution $P$ on $[-\pi, \pi]$ is defined as:
\begin{equation}
E_P[|\theta|] = \int_{-\pi}^\pi |\theta| p(\theta) d\theta = 2 \int_0^\pi \theta p(\theta) d\theta
\end{equation}
Using integration by parts with $u = \theta$ and $dv = p(\theta)d\theta$ (where $v = F_P(\theta) - 1$), we obtain:
\begin{align}
2 \int_0^\pi \theta p(\theta) d\theta &= 2 \left[ \theta (F_P(\theta) - 1) \right]_0^\pi - 2 \int_0^\pi (F_P(\theta) - 1) d\theta \nonumber \\
&= 2 \pi (1 - 1) - 0 + 2 \int_0^\pi (1 - F_P(\theta)) d\theta \nonumber \\
&= 2 \int_0^\pi (1 - F_P(\theta)) d\theta
\end{align}
Thus, we have:
\begin{equation}
E_P[|\theta|] = 2 \int_0^\pi (1 - F_P(\theta)) d\theta, \quad E_Q[|\theta|] = 2 \int_0^\pi (1 - F_Q(\theta)) d\theta
\end{equation}
Substituting these expressions back into the simplified 1-Wasserstein integral yields:
\begin{align}
W_1(P, Q) &= 2 \int_0^\pi (F_P(\theta) - F_Q(\theta)) d\theta \nonumber \\
&= 2 \int_0^\pi (1 - F_Q(\theta)) d\theta - 2 \int_0^\pi (1 - F_P(\theta)) d\theta \nonumber \\
&= E_Q[|\theta|] - E_P[|\theta|]
\end{align}
Taking the absolute value to handle both cases where either $P$ or $Q$ is more concentrated, we arrive at the desired result:
\begin{equation}
W_1(P, Q) = | E_P[|\theta|] - E_Q[|\theta|] |
\end{equation}

\subsection{Analytical Expressions for Symmetric Distributions}
The mean absolute deviation of each symmetric distribution is given by:
\begin{align}
E_P[|\theta|] &= \frac{\pi}{2} - \frac{4}{\pi} \sum_{k=0}^\infty \frac{I_{2k+1}(\kappa)}{(2k+1)^2 I_0(\kappa)} \\
E_Q[|\theta|] &= \frac{\pi}{2} - \frac{4}{\pi} \sum_{k=0}^\infty \frac{\rho^{2k+1}}{(2k+1)^2}
\end{align}
Thus, the 1-Wasserstein distance is:
\begin{equation}
W_1(P, Q) = \frac{4}{\pi} \left| \sum_{k=0}^\infty \frac{\rho^{2k+1} - \frac{I_{2k+1}(\kappa)}{I_0(\kappa)}}{(2k+1)^2} \right|
\end{equation}

\end{document}
